\newif\ifusectex
\IfFileExists{ctexart.cls}{%
  \usectextrue
  \documentclass[UTF8,12pt,a4paper]{ctexart}
}{%
  \usectexfalse
  \documentclass[12pt,a4paper]{article}
}

\ifusectex\else
  \usepackage{fontspec}
  \IfFileExists{output/pdf/fonts/NotoSansCJKsc-Regular.otf}{%
    \def\localcjkfontpath{output/pdf/fonts/}%
  }{%
    \def\localcjkfontpath{fonts/}%
  }
  \setmainfont[
    Path=\localcjkfontpath,
    AutoFakeBold=2,
    AutoFakeSlant=0.2
  ]{NotoSansCJKsc-Regular.otf}
  \setsansfont[
    Path=\localcjkfontpath,
    AutoFakeBold=2,
    AutoFakeSlant=0.2
  ]{NotoSansCJKsc-Regular.otf}
  \setmonofont{DejaVu Sans Mono}
  \XeTeXlinebreaklocale "zh"
  \XeTeXlinebreakskip=0pt plus 1pt
  \renewcommand{\contentsname}{目录}
\fi

\usepackage{amsmath,amssymb,mathtools,bm}
\usepackage{geometry}
\usepackage{booktabs,longtable,array}
\usepackage{enumitem}
\usepackage{xcolor}
\usepackage{fancyhdr}
\usepackage{hyperref}
\usepackage{listings}
\usepackage[most]{tcolorbox}

\geometry{left=2.35cm,right=2.35cm,top=2.45cm,bottom=2.45cm}
\setlength{\parindent}{2em}
\setlength{\parskip}{0.35em}
\linespread{1.25}
\allowdisplaybreaks[4]
\numberwithin{equation}{section}

\definecolor{mainblue}{HTML}{1F4E79}
\definecolor{softblue}{HTML}{EAF2F8}
\definecolor{sourcegray}{HTML}{555555}
\definecolor{codebg}{HTML}{F5F7FA}

\hypersetup{
  colorlinks=true,
  linkcolor=mainblue,
  urlcolor=mainblue,
  citecolor=mainblue,
  pdftitle={第三题：对流扩散方程的有限体积法完整推导},
  pdfauthor={张云飞}
}

\pagestyle{fancy}
\fancyhf{}
\fancyhead[L]{第三题：对流扩散方程}
\fancyhead[R]{有限体积法全显式求解器}
\fancyfoot[C]{\thepage}
\setlength{\headheight}{16pt}
\renewcommand{\headrulewidth}{0.4pt}

\ifusectex
  \ctexset{
    section={format=\Large\bfseries\color{mainblue},beforeskip=1.2ex,afterskip=0.8ex},
    subsection={format=\large\bfseries,beforeskip=1.0ex,afterskip=0.5ex},
    subsubsection={format=\normalsize\bfseries,beforeskip=0.8ex,afterskip=0.3ex}
  }
\else
  \usepackage{titlesec}
  \titleformat{\section}{\Large\bfseries\color{mainblue}}{\thesection}{0.7em}{}
  \titleformat{\subsection}{\large\bfseries}{\thesubsection}{0.7em}{}
  \titleformat{\subsubsection}{\normalsize\bfseries}{\thesubsubsection}{0.7em}{}
\fi

\newcommand{\vect}[1]{\bm{#1}}
\newcommand{\dd}{\mathop{}\!\mathrm{d}}
\newcommand{\src}[1]{\par\noindent{\small\color{sourcegray}\textbf{参考来源：}#1}\par}

\newtcolorbox{resultbox}{
  colback=softblue,
  colframe=mainblue,
  boxrule=0.7pt,
  arc=1.5mm,
  left=2mm,right=2mm,top=1.5mm,bottom=1.5mm
}

\lstset{
  basicstyle=\ttfamily\small,
  backgroundcolor=\color{codebg},
  frame=single,
  rulecolor=\color{gray!45},
  breaklines=true,
  columns=fullflexible,
  showstringspaces=false
}

\title{\bfseries 第三题：对流扩散方程\\[0.4em]
\Large 有限体积法半离散、全显式离散与求解算法}
\author{张云飞}
\date{2026年8月24日}

\begin{document}

\maketitle

\begin{tcolorbox}[colback=white,colframe=mainblue,title=离散方案]
本文严格按照题目给出的控制方程和此前提供的三份参考资料，采用如下方案：
\begin{itemize}[leftmargin=2.2em,itemsep=0.25em]
  \item 对流项：Gauss 有限体积积分与一阶迎风格式；
  \item 扩散项：Gauss 有限体积积分与中心差分；非正交网格采用 corrected 修正；
  \item 时间项：一阶前向欧拉格式；
  \item 显式性：所有空间通量均由旧时刻 $\phi^n$ 计算，不组装和求解隐式线性方程组。
\end{itemize}
\end{tcolorbox}

题目控制方程为
\begin{equation}
\boxed{
\frac{\partial\phi}{\partial t}
+\nabla\cdot(\vect u\phi)
-\nabla\cdot(\mu\nabla\phi)=0 .
}
\label{eq:governing}
\end{equation}
它对应参考书一般标量输运方程
\begin{equation}
\frac{\partial(\rho\phi)}{\partial t}
+\nabla\cdot(\rho\vect u\phi)
=\nabla\cdot(\Gamma^\phi\nabla\phi)+Q^\phi
\end{equation}
中的
\begin{equation}
\rho=1,\qquad \Gamma^\phi=\mu,\qquad Q^\phi=0.
\end{equation}
\src{R1，第5.2节，式(5.1)。}

\tableofcontents
\newpage

%============================================================
\section{第一部分：有限体积法半离散形式}

\subsection{控制体、单元平均值与面几何量}

将计算区域划分为互不重叠的控制体。对任意控制体 $P$，定义其体积
\begin{equation}
V_P=\int_{V_P}\dd V,
\end{equation}
以及单元平均值
\begin{equation}
\phi_P(t)=\frac{1}{V_P}\int_{V_P}\phi(\vect x,t)\dd V.
\label{eq:cell-average}
\end{equation}
控制体边界写成各个面的并：
\begin{equation}
\partial V_P=\bigcup_{f\in\partial P}f.
\end{equation}
对控制体 $P$ 的面 $f$，定义相对于 $P$ 向外的面积矢量
\begin{equation}
\vect S_{P,f}
=\int_f\vect n_P\dd S
\approx \vect n_{P,f}A_f,
\label{eq:area-vector}
\end{equation}
其中 $A_f$ 是面面积，$\vect n_{P,f}$ 是单位外法向量。对于内部面，若其 owner 单元为 $P$、neighbour 单元为 $N$，则
\begin{equation}
\vect S_{N,f}=-\vect S_{P,f}.
\label{eq:area-opposite}
\end{equation}

三维情况下，$\int_{V_P}\dd V$ 是三重体积分，$\int_{\partial V_P}\dd S$ 是二重曲面积分；题目中的二维算例则分别退化为面积积分和边界线积分。参考书采用 $\dd V$ 与 $\dd S$ 的坐标无关写法，因而不需要在推导中分别展开成直角坐标的二重或三重积分。
\src{R1，第5.2节，式(5.3)--(5.4)；R2，式(1.16)--(1.17)。}

\subsection{在控制体上积分}

将式\eqref{eq:governing}在控制体 $V_P$ 上积分：
\begin{equation}
\int_{V_P}\frac{\partial\phi}{\partial t}\dd V
+\int_{V_P}\nabla\cdot(\vect u\phi)\dd V
-\int_{V_P}\nabla\cdot(\mu\nabla\phi)\dd V=0.
\label{eq:volume-integral}
\end{equation}

控制体固定不动，因此
\begin{align}
\int_{V_P}\frac{\partial\phi}{\partial t}\dd V
&=\frac{\dd}{\dd t}\int_{V_P}\phi\dd V \\
&=V_P\frac{\dd\phi_P}{\dd t}.
\label{eq:transient-integral}
\end{align}

对对流项应用 Gauss 散度定理：
\begin{equation}
\int_{V_P}\nabla\cdot(\vect u\phi)\dd V
=\int_{\partial V_P}\vect u\phi\cdot\vect n_P\dd S.
\label{eq:gauss-advection}
\end{equation}
对扩散项同样应用 Gauss 散度定理：
\begin{equation}
\int_{V_P}\nabla\cdot(\mu\nabla\phi)\dd V
=\int_{\partial V_P}\mu\nabla\phi\cdot\vect n_P\dd S.
\label{eq:gauss-diffusion}
\end{equation}
代回式\eqref{eq:volume-integral}：
\begin{equation}
V_P\frac{\dd\phi_P}{\dd t}
+\int_{\partial V_P}\vect u\phi\cdot\vect n_P\dd S
-\int_{\partial V_P}\mu\nabla\phi\cdot\vect n_P\dd S=0.
\label{eq:integral-balance}
\end{equation}
\src{R1，第5.2节，式(5.3)--(5.10)。}

\subsection{由边界积分得到面通量求和}

把控制体边界积分分解成所有面的积分，并在每个面上使用单点 Gauss 积分：
\begin{align}
\int_f\vect u\phi\cdot\vect n_P\dd S
&\approx (\vect u\phi)_f\cdot\vect S_{P,f}, \\
\int_f\mu\nabla\phi\cdot\vect n_P\dd S
&\approx \mu_f(\nabla\phi)_f\cdot\vect S_{P,f}.
\end{align}
于是
\begin{equation}
V_P\frac{\dd\phi_P}{\dd t}
+\sum_{f\in\partial P}(\vect u\phi)_f\cdot\vect S_{P,f}
-\sum_{f\in\partial P}\mu_f(\nabla\phi)_f\cdot\vect S_{P,f}=0.
\label{eq:face-sum-before-F}
\end{equation}

定义相对于控制体 $P$ 向外的面体积通量
\begin{equation}
F_{P,f}=\vect u_f\cdot\vect S_{P,f}.
\label{eq:face-flux}
\end{equation}
于是
\begin{equation}
(\vect u\phi)_f\cdot\vect S_{P,f}=F_{P,f}\phi_f.
\end{equation}
因此有限体积半离散骨架为
\begin{resultbox}
\begin{equation}
V_P\frac{\dd\phi_P}{\dd t}
+\sum_{f\in\partial P}F_{P,f}\phi_f
-\sum_{f\in\partial P}\mu_f(\nabla\phi)_f\cdot\vect S_{P,f}=0.
\label{eq:semidiscrete-skeleton}
\end{equation}
等价地，
\begin{equation}
\frac{\dd\phi_P}{\dd t}
=-\frac{1}{V_P}\sum_{f\in\partial P}F_{P,f}\phi_f
+\frac{1}{V_P}\sum_{f\in\partial P}\mu_f(\nabla\phi)_f\cdot\vect S_{P,f}.
\label{eq:semidiscrete-rhs}
\end{equation}
\end{resultbox}

式\eqref{eq:semidiscrete-skeleton}已经完成“控制体积分、Gauss 定理、面通量求和”的过程；时间方向仍保持连续，因此称为半离散形式。但面值 $\phi_f$ 和面法向梯度 $(\nabla\phi)_f$ 尚未由单元值表示，故还需要空间离散使方程封闭。
\src{R1，第5.2节，式(5.8)--(5.13)；R2，式(1.22)--(1.23)。}

%============================================================
\section{第二部分：空间显式离散}

\subsection{对流项的一阶迎风离散}

考虑连接 owner 单元 $P$ 与 neighbour 单元 $N$ 的内部面 $f$。以下统一采用 owner 单元 $P$ 的外法向，简记
\begin{equation}
F_f=\vect u_f\cdot\vect S_{P,f}.
\end{equation}
若 $F_f\geq0$，流体从 $P$ 流向 $N$，上游单元为 $P$；若 $F_f<0$，流体从 $N$ 流向 $P$，上游单元为 $N$。所以
\begin{equation}
\phi_f^{\mathrm{up}}
=\begin{cases}
\phi_P,&F_f\geq0,\\[1mm]
\phi_N,&F_f<0.
\end{cases}
\label{eq:upwind-piecewise}
\end{equation}
定义
\begin{equation}
F_f^+=\max(F_f,0),\qquad F_f^-=\min(F_f,0),
\label{eq:Fplusminus}
\end{equation}
则
\begin{equation}
F_f=F_f^++F_f^-,
\end{equation}
且一阶迎风对流通量可统一写成
\begin{resultbox}
\begin{equation}
F_f\phi_f^{\mathrm{up}}
=F_f^+\phi_P+F_f^-\phi_N.
\label{eq:upwind-flux}
\end{equation}
\end{resultbox}
\src{R1，第11.2.4节，式(11.33)--(11.42)；R3，第59.2.1节。}

本题要求空间离散为显式，因此实际计算时所有量均取旧时刻 $n$：
\begin{equation}
\left(F_f\phi_f\right)^n
=(F_f^n)^+\phi_P^n+(F_f^n)^-\phi_N^n.
\label{eq:explicit-upwind}
\end{equation}
这里“显式”不仅指选择了迎风插值，还要求通量表达式中不存在待求的 $\phi^{n+1}$。

\subsection{正交网格上的中心扩散离散}

设
\begin{equation}
\vect d_{PN}=\vect x_N-\vect x_P,
\qquad d_{PN}=|\vect d_{PN}|.
\end{equation}
在正交网格上，单元中心连线与面法向平行，面法向梯度用中心差分近似：
\begin{equation}
(\nabla\phi)_f\cdot\vect n_f
\approx\frac{\phi_N-\phi_P}{d_{PN}}.
\end{equation}
故
\begin{equation}
(\nabla\phi)_f\cdot\vect S_{P,f}
\approx A_f\frac{\phi_N-\phi_P}{d_{PN}},
\end{equation}
从而
\begin{equation}
\mu_f(\nabla\phi)_f\cdot\vect S_{P,f}
\approx\frac{\mu_fA_f}{d_{PN}}(\phi_N-\phi_P).
\end{equation}
定义
\begin{equation}
D_f=\frac{\mu_fA_f}{d_{PN}},
\label{eq:D-orthogonal}
\end{equation}
正交网格上的扩散通量为
\begin{resultbox}
\begin{equation}
G_f:=\mu_f(\nabla\phi)_f\cdot\vect S_{P,f}
\approx D_f(\phi_N-\phi_P).
\label{eq:diffusion-orthogonal}
\end{equation}
\end{resultbox}
\src{R1，第8.1节，式(8.9)--(8.21)。}

\subsection{非正交网格的 corrected 扩散离散}

题目要求在三角形和四边形网格上计算，实际网格一般并不完全正交。按照参考书，将面面积矢量分解为
\begin{equation}
\vect S_{P,f}=\vect E_f+\vect T_f,
\label{eq:nonorth-decomp}
\end{equation}
其中 $\vect E_f$ 与 $\vect d_{PN}$ 平行，$\vect T_f$ 表示非正交部分。于是
\begin{align}
(\nabla\phi)_f\cdot\vect S_{P,f}
&=(\nabla\phi)_f\cdot\vect E_f
+(\nabla\phi)_f\cdot\vect T_f \\
&\approx\frac{|\vect E_f|}{d_{PN}}(\phi_N-\phi_P)
+(\nabla\phi)_f\cdot\vect T_f.
\end{align}
乘以 $\mu_f$，定义
\begin{equation}
D_f=\mu_f\frac{|\vect E_f|}{d_{PN}},
\qquad
C_f=\mu_f(\nabla\phi)_f\cdot\vect T_f,
\end{equation}
得到
\begin{resultbox}
\begin{equation}
G_f
=D_f(\phi_N-\phi_P)+C_f.
\label{eq:diffusion-corrected}
\end{equation}
\end{resultbox}
\src{R1，第8.6节，式(8.63)--(8.80)。}

为使修正项仍然显式，使用旧时刻梯度：
\begin{equation}
C_f^n=\mu_f(\nabla\phi)_f^n\cdot\vect T_f.
\end{equation}
单元梯度按照 Gauss 梯度公式计算：
\begin{equation}
(\nabla\phi)_P^n
=\frac{1}{V_P}\sum_{f\in\partial P}\phi_f^n\vect S_{P,f}.
\label{eq:gauss-gradient}
\end{equation}
再将相邻单元梯度几何插值到面中心：
\begin{equation}
(\nabla\phi)_f^n
=\lambda_f(\nabla\phi)_P^n
+(1-\lambda_f)(\nabla\phi)_N^n.
\label{eq:face-gradient-interpolation}
\end{equation}
\src{R1，第8.6节，式(8.75)--(8.76)。}

\subsection{合并后的封闭空间半离散方程}

将一阶迎风对流通量\eqref{eq:upwind-flux}与扩散通量\eqref{eq:diffusion-corrected}代入式\eqref{eq:semidiscrete-skeleton}：
\begin{align}
V_P\frac{\dd\phi_P}{\dd t}
=&-\sum_{f\in\partial P}
\left(F_f^+\phi_P+F_f^-\phi_N\right) \\
&+\sum_{f\in\partial P}
\left[D_f(\phi_N-\phi_P)+C_f\right].
\label{eq:closed-semidiscrete}
\end{align}
整理中心单元和相邻单元系数：
\begin{resultbox}
\begin{align}
V_P\frac{\dd\phi_P}{\dd t}
=&-\left[\sum_{f\in\partial P}(F_f^++D_f)\right]\phi_P \\
&+\sum_{f\in\partial P}(D_f-F_f^-)\phi_N
+\sum_{f\in\partial P}C_f.
\label{eq:closed-coefficient}
\end{align}
\end{resultbox}
因为
\begin{equation}
F_f^+\geq0,\qquad F_f^-\leq0,\qquad D_f\geq0,
\end{equation}
所以相邻单元系数 $D_f-F_f^-\geq0$。

\subsection{均匀直角网格上的具体展开}

设二维均匀直角网格步长为 $\Delta x,\Delta y$，速度和扩散系数为常数：
\begin{equation}
\vect u=(u,v),\qquad \mu=\text{const}.
\end{equation}
定义
\begin{equation}
u^+=\max(u,0),\quad u^-=\min(u,0),\quad
v^+=\max(v,0),\quad v^-=\min(v,0).
\end{equation}
则式\eqref{eq:closed-semidiscrete}退化为
\begin{align}
\frac{\dd\phi_{i,j}}{\dd t}
=&-\frac{u^+(\phi_{i,j}-\phi_{i-1,j})
+u^-(\phi_{i+1,j}-\phi_{i,j})}{\Delta x} \\
&-\frac{v^+(\phi_{i,j}-\phi_{i,j-1})
+v^-(\phi_{i,j+1}-\phi_{i,j})}{\Delta y} \\
&+\mu\frac{\phi_{i+1,j}-2\phi_{i,j}+\phi_{i-1,j}}{\Delta x^2} \\
&+\mu\frac{\phi_{i,j+1}-2\phi_{i,j}+\phi_{i,j-1}}{\Delta y^2}.
\label{eq:cartesian-general}
\end{align}
若 $u>0,v>0$，则
\begin{resultbox}
\begin{align}
\frac{\dd\phi_{i,j}}{\dd t}
=&-u\frac{\phi_{i,j}-\phi_{i-1,j}}{\Delta x}
-v\frac{\phi_{i,j}-\phi_{i,j-1}}{\Delta y} \\
&+\mu\frac{\phi_{i+1,j}-2\phi_{i,j}+\phi_{i-1,j}}{\Delta x^2} \\
&+\mu\frac{\phi_{i,j+1}-2\phi_{i,j}+\phi_{i,j-1}}{\Delta y^2}.
\label{eq:cartesian-positive}
\end{align}
\end{resultbox}

%============================================================
\section{第三部分：前向欧拉时间离散与全显式求解器}

\subsection{前向欧拉格式的推导}

在 $t^n$ 对 $\phi_P(t)$ 作 Taylor 展开：
\begin{equation}
\phi_P(t^{n+1})
=\phi_P(t^n)
+\Delta t\left.\frac{\dd\phi_P}{\dd t}\right|_{t^n}
+O(\Delta t^2).
\end{equation}
移项并除以 $\Delta t$：
\begin{equation}
\left.\frac{\dd\phi_P}{\dd t}\right|_{t^n}
=\frac{\phi_P^{n+1}-\phi_P^n}{\Delta t}+O(\Delta t).
\end{equation}
忽略截断误差，得到一阶前向欧拉格式
\begin{resultbox}
\begin{equation}
\left.\frac{\dd\phi_P}{\dd t}\right|_{t^n}
\approx\frac{\phi_P^{n+1}-\phi_P^n}{\Delta t}.
\label{eq:forward-euler}
\end{equation}
\end{resultbox}
\src{R1，第13.2.1节，式(13.5)--(13.12)。}

\subsection{完全显式更新公式}

将式\eqref{eq:forward-euler}代入式\eqref{eq:closed-semidiscrete}，并把所有空间量放在旧时刻 $n$ 计算：
\begin{align}
V_P\frac{\phi_P^{n+1}-\phi_P^n}{\Delta t}
=&-\sum_{f\in\partial P}
\left[(F_f^n)^+\phi_P^n+(F_f^n)^-\phi_N^n\right] \\
&+\sum_{f\in\partial P}
\left[D_f^n(\phi_N^n-\phi_P^n)+C_f^n\right].
\label{eq:fully-discrete-before-update}
\end{align}
因此
\begin{resultbox}
\begin{align}
\phi_P^{n+1}
=\phi_P^n+\frac{\Delta t}{V_P}\Bigg\{
&-\sum_{f\in\partial P}
\left[(F_f^n)^+\phi_P^n+(F_f^n)^-\phi_N^n\right] \\
&+\sum_{f\in\partial P}
\left[D_f^n(\phi_N^n-\phi_P^n)+C_f^n\right]
\Bigg\}.
\label{eq:fully-explicit-update}
\end{align}
\end{resultbox}

进一步展开中心单元和相邻单元系数：
\begin{align}
\phi_P^{n+1}
=&\left[1-\frac{\Delta t}{V_P}
\sum_{f\in\partial P}(F_f^++D_f)\right]\phi_P^n \\
&+\frac{\Delta t}{V_P}
\sum_{f\in\partial P}(D_f-F_f^-)\phi_N^n
+\frac{\Delta t}{V_P}\sum_{f\in\partial P}C_f^n.
\label{eq:explicit-coefficient-form}
\end{align}
式\eqref{eq:fully-explicit-update}或式\eqref{eq:explicit-coefficient-form}就是本题要求的全显式有限体积求解器。右端只含已知的旧时刻变量，不需要求解全局线性方程组。

\subsection{显式格式的稳定时间步}

按照参考书对多维显式对流--扩散方程的稳定性分析，时间步需要满足
\begin{resultbox}
\begin{equation}
\Delta t\leq
\min_P
\frac{V_P}{
\displaystyle\sum_{f\in\partial P}
\left[D_f+\max(F_f,0)\right]}.
\label{eq:general-dt}
\end{equation}
\end{resultbox}
实际计算引入安全系数 $0<\sigma<1$：
\begin{equation}
\Delta t
=\sigma\min_P
\frac{V_P}{
\displaystyle\sum_{f\in\partial P}(D_f+F_f^+)}.
\label{eq:safe-dt}
\end{equation}
可先取 $\sigma=0.5\sim0.8$。非正交修正项显式处理时，式\eqref{eq:general-dt}没有单独计入修正项的额外影响，因此高非正交网格应使用更保守的 $\sigma$。
\src{R1，第13.2.2.3节，式(13.25)--(13.29)。}

对二维均匀直角网格，式\eqref{eq:general-dt}化为
\begin{equation}
\Delta t\left[
\frac{|u|}{\Delta x}+\frac{|v|}{\Delta y}
+2\mu\left(\frac{1}{\Delta x^2}+\frac{1}{\Delta y^2}\right)
\right]\leq1,
\label{eq:cartesian-stability}
\end{equation}
即
\begin{equation}
\Delta t\leq
\frac{1}{
\dfrac{|u|}{\Delta x}+\dfrac{|v|}{\Delta y}
+2\mu\left(\dfrac{1}{\Delta x^2}+\dfrac{1}{\Delta y^2}\right)}.
\end{equation}
第一个算例中 $\vect u=(a,a)$、$\mu=1$ 且 $\Delta x=\Delta y=h$，故
\begin{equation}
\Delta t\leq\frac{1}{2|a|/h+4/h^2}.
\label{eq:sine-dt}
\end{equation}
题面没有给出 $a$ 的具体数值，实际计算前必须指定 $a$，否则不能确定数值时间步及精确解的具体平移距离。

%============================================================
\section{第四部分：总体算法、OpenFOAM 对应与算例验证}

\subsection{网格预处理与初始化}

读取网格，建立每个面的 owner--neighbour 关系，并计算
\begin{equation}
V_P,\quad \vect x_P,\quad \vect x_f,\quad
\vect S_{P,f},\quad d_{PN},\quad \vect E_f,\quad \vect T_f.
\end{equation}
计算扩散面系数
\begin{equation}
D_f=\mu_f\frac{|\vect E_f|}{d_{PN}}.
\end{equation}
初始化单元平均值
\begin{equation}
\phi_P^0=\frac{1}{V_P}\int_{V_P}\phi(\vect x,t_0)\dd V.
\label{eq:initial-cell-average}
\end{equation}
程序中可以采用单元中心近似 $\phi_P^0\approx\phi(\vect x_P,t_0)$；但若要严格研究网格收敛阶，宜使用单元积分或足够高阶的数值积分计算初始单元平均值。

\subsection{单个时间步的守恒通量装配}

在时刻 $t^n$，首先在面中心计算或插值速度：
\begin{equation}
\vect u_f^n=\vect u(\vect x_f,t^n),
\qquad F_f^n=\vect u_f^n\cdot\vect S_{P,f}.
\end{equation}
根据式\eqref{eq:safe-dt}计算 $\Delta t^n$，并在最后一步使用
\begin{equation}
\Delta t^n=\min(\Delta t^n,T-t^n)
\end{equation}
以保证恰好到达目标时间。

对于内部面 $f$，设其 owner 为 $P$、neighbour 为 $N$。迎风面值为
\begin{equation}
\phi_f^{\mathrm{up},n}
=\begin{cases}
\phi_P^n,&F_f^n\geq0,\\
\phi_N^n,&F_f^n<0.
\end{cases}
\end{equation}
定义对流通量、扩散通量和方程中的总向外通量：
\begin{align}
A_f^n&=F_f^n\phi_f^{\mathrm{up},n},\\
G_f^n&=D_f(\phi_N^n-\phi_P^n)+C_f^n,\\
J_f^n&=A_f^n-G_f^n.
\label{eq:combined-face-flux}
\end{align}
初始化 $R_P^n=0$，然后对每个内部面执行
\begin{resultbox}
\begin{equation}
R_P^n\mathrel{-}=J_f^n,
\qquad
R_N^n\mathrel{+}=J_f^n.
\label{eq:owner-neighbour-assembly}
\end{equation}
\end{resultbox}
同一面从 owner 流出的通量恰好等于 neighbour 流入的通量，故该装配方式保证内部面通量严格守恒。
\src{R2，式(1.24)--(1.26)及 owner--neighbour 面通量装配过程。}

所有面处理完后，使用
\begin{equation}
\phi_P^{n+1}=\phi_P^n+\frac{\Delta t^n}{V_P}R_P^n
\label{eq:residual-update}
\end{equation}
更新单元值。

\subsection{边界条件离散}

\subsubsection{周期边界}
将周期匹配面另一侧的单元视为 neighbour，按内部面通量装配。第一个正弦波算例的四周均使用周期边界。

\subsubsection{Dirichlet 边界}
若边界给定 $\phi=\phi_b$，对流通量按照迎风方向取值：
\begin{equation}
\phi_f^{\mathrm{up}}
=\begin{cases}
\phi_P,&F_b\geq0\quad\text{（流出）},\\
\phi_b,&F_b<0\quad\text{（流入）}.
\end{cases}
\end{equation}
扩散边界通量可近似为
\begin{equation}
G_b=D_b(\phi_b-\phi_P)+C_b,
\qquad D_b=\frac{\mu_bA_b}{d_{Pb}}.
\end{equation}

\subsubsection{Neumann 边界}
若给定 $\nabla\phi\cdot\vect n=g_b$，则扩散通量直接为
\begin{equation}
G_b=\mu_bg_bA_b.
\end{equation}
\src{R1，第8.3节，式(8.30)--(8.42)。}

\subsection{总体算法}

完整算法可写为：
\begin{enumerate}[leftmargin=2.5em,label=\arabic*.,itemsep=0.45em]
  \item 读取网格并计算单元、面和 owner--neighbour 几何信息；
  \item 根据精确初值初始化 $\phi_P^0$，令 $t=t_0$；
  \item 在所有面中心计算 $\vect u_f$ 和 $F_f=\vect u_f\cdot\vect S_f$；
  \item 根据式\eqref{eq:safe-dt}计算稳定时间步，并截断最后一个时间步；
  \item 由式\eqref{eq:gauss-gradient}--\eqref{eq:face-gradient-interpolation}计算旧时刻梯度及显式非正交修正 $C_f^n$；
  \item 将所有单元残差初始化为零；
  \item 遍历内部面，计算 $A_f^n,G_f^n,J_f^n$，按式\eqref{eq:owner-neighbour-assembly}守恒装配；
  \item 处理周期、Dirichlet 或 Neumann 边界面；
  \item 按式\eqref{eq:residual-update}作前向欧拉更新并修正边界值；
  \item 更新时间，输出所需时刻的数值场；若 $t<T$，返回步骤3；
  \item 在目标时刻计算误差范数和网格收敛阶。
\end{enumerate}

等价伪代码如下：
\begin{lstlisting}
readMesh();
precomputeGeometry();
initializePhi();

while (time < endTime)
{
    computeFaceVelocityAndFlux();
    dt = computeStableTimeStep();
    dt = min(dt, endTime - time);

    computeCellGradients(phiOld);
    computeExplicitNonOrthogonalCorrection(phiOld);

    residual = 0;
    forAllInternalFaces(f)
    {
        P = owner[f];
        N = neighbour[f];
        phiUp = (faceFlux[f] >= 0) ? phiOld[P] : phiOld[N];

        advFlux  = faceFlux[f]*phiUp;
        diffFlux = D[f]*(phiOld[N] - phiOld[P]) + correction[f];
        totalFlux = advFlux - diffFlux;

        residual[P] -= totalFlux;
        residual[N] += totalFlux;
    }

    applyBoundaryFluxes();
    phiNew = phiOld + dt*residual/cellVolume;
    correctBoundaryConditions();

    phiOld = phiNew;
    time += dt;
}
\end{lstlisting}

\subsection{OpenFOAM 显式算子的对应关系}

参考书明确区分：\texttt{fvc::} 为显式有限体积算子，返回已经计算出的场；\texttt{fvm::} 为隐式算子，返回待求线性系统的离散矩阵。因此，本题的全显式空间算子可写成
\begin{lstlisting}[language=C++]
volScalarField rhs
(
    -fvc::div(faceFlux, phi)
    +fvc::laplacian(mu, phi)
);

phi.primitiveFieldRef()
    += runTime.deltaTValue()*rhs.primitiveField();

phi.correctBoundaryConditions();
\end{lstlisting}
其中：
\begin{itemize}[leftmargin=2.2em,itemsep=0.2em]
  \item \texttt{-fvc::div(faceFlux, phi)} 对应负对流散度；
  \item \texttt{fvc::laplacian(mu, phi)} 对应正扩散贡献。
\end{itemize}
离散格式设置为
\begin{lstlisting}
gradSchemes
{
    default                 Gauss linear;
}

divSchemes
{
    default                 none;
    div(faceFlux,phi)       Gauss upwind;
}

laplacianSchemes
{
    default                 none;
    laplacian(mu,phi)       Gauss linear corrected;
}
\end{lstlisting}
\src{R1，第5.10.2节表5.1、第11.9.2节和第8.10.2节。}

\subsection{算例一：正弦波的验证}

题目给出
\begin{equation}
\phi(x,y,0)=\sin\bigl(2\pi(x+y)\bigr),
\end{equation}
\begin{equation}
\vect u=(a,a),\qquad \mu=1,
\end{equation}
计算域为 $[0,1]^2$，边界条件为周期边界，目标时间为 $T=1$。精确解为
\begin{equation}
\phi_{\mathrm{exact}}(x,y,t)
=\exp(-8\pi^2\mu t)
\sin\bigl(2\pi(x+y-2at)\bigr).
\label{eq:sine-exact}
\end{equation}

分别在逐步加密的三角形和四边形网格上计算体积加权误差：
\begin{align}
L_1&=\frac{\sum_PV_P|e_P|}{\sum_PV_P},\\
L_2&=\left(\frac{\sum_PV_Pe_P^2}{\sum_PV_P}\right)^{1/2},\\
L_\infty&=\max_P|e_P|,
\end{align}
其中
\begin{equation}
e_P=\phi_P-\phi_{\mathrm{exact},P}.
\end{equation}
相邻两层网格的观测收敛阶为
\begin{equation}
p=\frac{\ln(E_h/E_{h/2})}{\ln2}.
\label{eq:observed-order}
\end{equation}
一阶迎风对流离散具有一阶数值耗散；虽然中心扩散项在规则网格上为二阶，但整体空间误差通常由迎风项主导，因此平滑正弦波算例的观测空间收敛阶预期接近一阶。另一方面，扩散稳定性使 $\Delta t=O(h^2)$，故前向欧拉的 $O(\Delta t)$ 时间误差通常为 $O(h^2)$，不会掩盖一阶迎风空间误差。
\src{题目第3.2.1节给出初值、周期边界、精确解、网格类型和目标时间；迎风格式阶数依据 R1 第11.2.4节。}

\subsection{算例二：旋转尖峰的验证}

题目给出
\begin{equation}
\vect u=(-y,x),\qquad \mu=\varepsilon=10^{-3},
\end{equation}
计算域为 $[-1,1]^2$。每个面上的速度和通量直接由面中心坐标计算：
\begin{equation}
\vect u_f=(-y_f,x_f),
\qquad F_f=(-y_f,x_f)\cdot\vect S_f.
\end{equation}
题面解析函数为
\begin{equation}
\phi(\vect x,t)
=\frac{1}{4\pi\varepsilon t}
\exp\left(-\frac{r^2}{4\varepsilon t}\right),
\end{equation}
\begin{equation}
r^2=(x-\widehat x)^2+(y-\widehat y)^2.
\end{equation}
题面写出的峰中心为
\begin{equation}
\widehat x(t)=x_0\cos t-y_0\sin t,
\qquad
\widehat y(t)=y_0\cos t-x_0\sin t.
\label{eq:printed-trajectory}
\end{equation}
从 $t_0=\pi/2$ 开始，一圈旋转的时间长度为 $2\pi$，若按角速度为1理解，则终止时间为
\begin{equation}
T=t_0+2\pi=\frac{5\pi}{2}.
\end{equation}
在终止时刻绘制等值线，并使用前述 $L_1,L_2,L_\infty$ 及式\eqref{eq:observed-order}比较三角形、四边形网格的误差和收敛行为。

\begin{tcolorbox}[colback=yellow!7,colframe=orange!70!black,title=题面符号需在编码前确认]
速度场 $\dot x=-y,\ \dot y=x$ 对应的一般旋转轨迹应为
\[
\widehat x=x_0\cos t-y_0\sin t,
\qquad
\widehat y=x_0\sin t+y_0\cos t.
\]
题面式\eqref{eq:printed-trajectory}第二式中 $x_0\sin t$ 的符号与标准旋转解不同。不过本题指定 $x_0=0$，因此该符号差异在当前参数下不影响 $\widehat y=y_0\cos t$。此外，“从 $t_0=\pi/2$ 开始”与“峰位于 $(x_0,y_0)$”之间存在时间原点歧义：必须确认 $(x_0,y_0)$ 是 $t=0$ 还是 $t=t_0$ 的峰中心，不能在程序中混用两种定义。
\end{tcolorbox}

由于一阶迎风格式具有明显数值扩散，尖峰在计算中会同时受到物理扩散 $\varepsilon$ 和迎风数值扩散的展宽作用。因此，验证时除误差范数外，还应比较峰值高度、峰中心位置与等值线宽度；网格加密后数值扩散应逐渐减弱。
\src{题目第3.2.2节给出区域、速度场、扩散系数、解析函数、起始时间、峰位置及“一圈旋转”验证要求。}

\section*{参考资料}
\addcontentsline{toc}{section}{参考资料}

\begin{description}[leftmargin=3.2em,labelwidth=2.5em,style=nextline]
  \item[R1] F. Moukalled, L. Mangani, M. Darwish, \textit{The Finite Volume Method in Computational Fluid Dynamics: An Advanced Introduction with OpenFOAM and Matlab}, Springer, 2016.
  \item[R2] T. Holzmann, K. Rooney, \textit{The OpenFOAM Technology Primer}, First Edition, 2021.
  \item[R3] G. Holzinger, \textit{OpenFOAM: A Little User-Manual}, 2020.
  \item[题目] \textit{Training Examples}, 第3题 ``Advection and diffusion equation''，包括第3.1节控制方程与第3.2节两个验证算例。
\end{description}

\end{document}
